\documentclass[10pt,a4paper]{amsart}
\usepackage[margin=19mm]{geometry}
\usepackage{amsmath,amssymb,mathtools}
\usepackage[scr=rsfs,cal=euler]{mathalfa}
\usepackage{xfrac}
\usepackage[unicode,hidelinks,bookmarks=false]{hyperref}
\newcommand{\ie}{i.e.\ }
\newcommand{\cf}{cf.\ }
\newcommand{\resp}{respectively\ }
\newcommand{\Subsection}{\subsection}
\providecommand{\nobreakdash}{\nobreak\hbox{-}}
\setlength{\emergencystretch}{2em}
\multlinegap0pt
\usepackage{bm}
\usepackage{bbm}
\usepackage{dsfont}
\usepackage{xspace}
\usepackage{cancel}
\usepackage{mathtools}
\usepackage{amsthm}
\makeatletter
\@ifundefined{subsection}{\newtheorem{subsection}{Subsection}}{}
\@ifundefined{theorem}{\newtheorem{theorem}[subsection]{Theorem}}{}
\makeatother
\newcommand{\D}[1]{\mathbb{#1}}% Doubled - blackboard bold - only caps, uas as \D{A}
\newcommand{\UU}{\bigcup}% big union
\newcommand{\ci}{\subseteq}% contained in with equality
\newcommand{\ga}{\alpha}
\newcommand{\gb}{\beta}
\newcommand{\set}[1]{\{#1\}}% set
\newcommand{\te}{\text}% same as \mathrm command.
\newcommand{\tit}{\textit}% text italic
\newcommand{\uu}{\cup}% union
\renewcommand{\gg}{\gamma}% old use >>
\title[Conditional constrained and unconstrained quantization for p]{Conditional constrained and unconstrained quantization for probability distributions}
\author{Megha Pandey, Mrinal Kanti Roychowdhury}
\date{Source: arXiv:2312.02965v3 (CC BY 4.0)}
\begin{document}
\maketitle

\noindent\emph{Source and licence.} Conditional constrained and unconstrained quantization for probability distributions. Version arXiv:2312.02965v3 (CC BY 4.0), distributed under the Creative Commons Attribution 4.0 International licence, as recorded in the arXiv licence metadata for this exact version and shown on the abstract page https://arxiv.org/abs/2312.02965v3. Full rights notice: licenses/arxiv-2312.02965-CC-BY-4.0.txt.\par\smallskip

\noindent\textit{Editorial scope.} Counted unit: the Theorem whose source label is \texttt{TheoM2} in the article (source0.tex lines 1070--1072, source label \texttt{TheoM2}), delivered together with the article's own complete original proof of it (source0.tex lines 1075--1108, one \texttt{proof} environment of 34 lines). No mathematics is written, repaired or re-derived: statement and proof are the article's own text line for line. Required context retained uncounted: none. Every definition, display and label of those lines is retained unchanged. Omitted from this package and not redistributed: 1--1069, 1073--1074, 1109--1230. Nothing inside a retained excerpt is omitted or altered: every retained body is delivered exactly as the article's lines are written, so no omission receipt of any kind accompanies this package. Citations: the retained excerpts carry no citation command at all, so no bibliography entry of the article is transcribed and no reference is redistributed. Reference adaptation: none was needed and none was made; every reference inside the delivered unit resolves inside it, so no unresolved reference or citation is produced. Printed slips kept literal: no printed word, symbol, hypothesis or number of the counted statement or of its proof was altered. Layout and class: the article's own class is \texttt{amsart} with packages including amssymb, latexsym, amsmath, amscd, amsthm, graphicx, color, xy, pgf, tikz, mathrsfs, geometry, hyperref; the wrapper is the lane's pinned \texttt{amsart} editorial wrapper at 10pt on A4, which restates the theorem-like environments the retained text needs and adds the article's own macro definitions that the retained text needs (\texttt{\textbackslash D}, \texttt{\textbackslash UU}, \texttt{\textbackslash ci}, \texttt{\textbackslash ga}, \texttt{\textbackslash gb}, \texttt{\textbackslash gg}, \texttt{\textbackslash set}, \texttt{\textbackslash te}, \texttt{\textbackslash tit}, \texttt{\textbackslash uu}), with extra wrapper packages (bm, bbm, dsfont, xspace, cancel, mathtools). The delivered numbering is the unit's own. Assets: no retained span contains a figure environment, a graphics-inclusion command, a PDF-page inclusion, a table or a TikZ picture; no image file is copied into this package, the package declares no asset dependency and no image file is redistributed with it. Licence: version arXiv:2312.02965v3 of this article is distributed under the Creative Commons Attribution 4.0 International licence, as recorded in the arXiv licence metadata for this exact version and shown on the abstract page https://arxiv.org/abs/2312.02965v3. A full rights notice is delivered at licenses/arxiv-2312.02965-CC-BY-4.0.txt. Characters: every retained excerpt is pure ASCII, so the delivered engine, XeLaTeX with its default Latin Modern text font, reports no missing character.
\begin{theorem} \label{TheoM2} 
Let $P$ be a Borel probability measure on $\D R^k$. In conditional constrained quantization, let the conditional set is contained in the union of the family of constraints. Then, in the conditional constrained quantization, the lower and upper quantization dimensions, and the lower and upper quantization coefficients for a Borel probability measure do not depend on the conditional set. 
\end{theorem} 

\begin{proof} 
In conditional constrained quantization, let $\set{S_j\ci \D R^k: j\in \D N}$ be the family of constraints, and let the set $\gb$ be the conditional set with $\te{card}(\gb)=\ell$ for some $\ell\in \D N$ such that the following condition is true: 
\begin{equation*}
\gb \ci \UU_{j\in \D N} S_j.
\end{equation*}
Let $q$ be the least positive integer such that $\gb\ci \mathop{\uu}\limits_{j=1}^q S_j$.  
Let  $ V_{n,r}(P)$ and  $V_{c,n,r}(P)$ denote  the $n$th constrained  and the $n$th conditional constrained quantization  errors, respectively. 
We now prove the following claim.

\tit{Claim.  Let $n> \max \set{q, \ell}$. Then, 
$V_{n,r}(P)\leq V_{c, n,r}(P)\leq V_{n-\ell,r}(P).$}

The proof of $V_{n,r}(P)\leq V_{c, n,r}(P)$ is trivial. We now give the proof of $V_{c, n,r}(P)\leq V_{n-\ell,r}(P)$. 
Take any $\gg\ci \mathop{\uu}\limits_{j=1}^q S_j$ with $1\leq \te{card}(\gg)\leq n-\ell$. Then, for any $x\in \D R^k$, we have 
\[\min_{a\in \gg\uu \gb}d(x, a)^r \leq \min_{a\in \gg} d(x, a)^r.\]
Hence, 
\begin{align*} 
V_{c, n, r}(P)&=\inf_{\ga} \Big\{\int \mathop{\min}\limits_{a\in\ga\uu\gb} d(x, a)^r dP(x) : \ga \ci \mathop{\uu}\limits_{j=1}^q S_j, ~ 1\leq  \text{card}(\ga) \leq n-\ell \Big\}\\
&\leq \int \mathop{\min}\limits_{a\in\gg\uu\gb} d(x, a)^r dP(x)\\
&\leq \int \mathop{\min}\limits_{a\in\gg} d(x, a)^r dP(x),
\end{align*} 
which yields the fact that 
\[V_{c, n, r}(P)\leq \inf_{\gg} \Big\{\int \mathop{\min}\limits_{a\in\gg } d(x, a)^r dP(x) : \gg\ci \mathop{\uu}\limits_{j=1}^q S_j, ~ 1\leq  \text{card}(\gg) \leq n-\ell \Big\}=V_{n-\ell, r}(P).\]
Thus, the proof of the claim is complete. 

By the claim, using the squeeze theorem, we have $V_{\infty, r}(P)=V_{c, \infty, r}(P)$. Hence, 
\begin{align}\label{eq8}
	V_{n,r}(P)-V_{\infty, r}(P)\leq V_{c, n,r}(P)-V_{c, \infty, r}(P) \leq V_{n-\ell,r}(P)-V_{\infty, r}(P).
\end{align}
By the inequalities in \eqref{eq8}, and applying the squeeze theorem, we conclude that the lower and upper quantization coefficients are identical in both the constrained and the conditionally constrained cases. 
 Choose $n\in \mathbb{N}$ large enough such that $V_{n-\ell,r}(P)-V_{\infty, r}(P)<1$. Then, by the inequalities in \eqref{eq8}, we get  
 $$\frac{r\log n}{-\log(V_{n,r}(P)-V_{\infty, r}(P))}\leq \frac{r\log n}{-\log(V_{c, n,r}(P)-V_{c, \infty, r}(P))}\leq \frac{r\log n}{-\log(V_{n-\ell,r}(P)-V_{\infty, r}(P))}.$$
Hence, by the squeeze theorem, it follows that the lower and upper quantization dimensions are equal in both the constrained and the conditionally constrained settings. This completes the proof of the theorem.
\end{proof} 

\end{document}
